\fchapter{Ej 11 Pág 84}

\fsection{\boldit{\textcolor{waveBlue2}{Actividad}}} 
\subsection*{Además de los organismos públicos, las empresas privadas llevan a cabo iniciativas relacionadas con los océanos y mares. En grupos de cuatro visitad la página web \href{www.iberdrola.com/sostenibilidad/limpieza-oceanos-iniciativas}{iberdrola}  y realizad una pequeña presentación sobre las iniciativas de IBERDROLA, entablando un posterior debate sobre cómo podéis contribuir a la consecución del ODS 14 en vuestro ámbito profesional.}

\begin{solution}
IBERDROLA contribuye al ODS 14 (Vida Submarina) a través de:
\begin{enumerate}
\item Voluntariado ambiental: Organiza limpiezas de playas y fondos marinos.
\item Colaboraciones: Trabaja con ONGs como SEO/BirdLife en proyectos de conservación.
\item Investigación: Financia estudios científicos sobre ecosistemas marinos.
\item Estrategia: Incluye la "economía azul" en su plan de sostenibilidad.
\end{enumerate}
En nuestro ámbito profesional, podemos contribuir de las siguientes maneras:
\begin{enumerate}
	\item \textbf{EN PROYECTOS COSTEROS O INDUSTRIALES:}
\begin{itemize}
	\item{Instala sensores y sistemas automáticos para contener vertidos en puertos o industrias.}
	\item{Automatiza plantas de tratamiento de aguas para que no lleguen contaminantes al mar.}
	\item{Diseña sistemas de energías renovables (fotovoltaica) para alimentar instalaciones costeras.}
\end{itemize}

\item \textbf{EN ECONOMÍA AZUL SOSTENIBLE:}

\begin{itemize}
	\item{Automatiza granjas marinas para un uso eficiente de comida y control de la calidad del agua.}
	\item{Implementa sistemas de electricidad en puerto para barcos ("cold ironing") para reducir emisiones.}
	\item{Mantén redes de sensores marinos para investigación oceanográfica.}
\end{itemize}

	\item \textbf{EN EL TRABAJO DIARIO:}

\begin{itemize}
	\item{Gestiona correctamente los residuos de obra (cables, baterías, RAEE) para que no contaminen.}
	\item{Reduce el uso de plásticos y materiales contaminantes en las instalaciones.}
\end{itemize}

\end{enumerate}
\end{solution} 



